\section{Introduction and Background}

A bright and statistically significant excess of GeV-scale gamma-ray emission has been observed from the region surrounding the Galactic Center in the data collected by the Fermi Gamma-Ray Space Telescope~\cite{Goodenough:2009gk,Hooper:2010mq,Hooper:2011ti,Abazajian:2012pn,Hooper:2013rwa,Gordon:2013vta,Daylan:2014rsa,Calore:2014xka,Fermi-LAT:2015sau,Fermi-LAT:2017opo}. This Galactic Center Gamma-Ray Excess (GCE) persists across all existing models of the astrophysical diffuse emission and has not been attributed to any known astrophysical sources or mechanisms~\cite{Cholis:2021rpp,DiMauro:2021raz}. The observed spectrum, angular distribution, and overall intensity of this excess are each consistent with the signal predicted from the annihilation of dark matter particles in the form of a $m\sim 50 \, \textrm{GeV}$ thermal relic.

The most prominent astrophysical interpretation of the GCE is that it could be generated by a large population of centrally-located millisecond pulsars (MSPs)~\cite{Hooper:2010mq,Abazajian:2010zy,Hooper:2011ti,Abazajian:2012pn,Gordon:2013vta,Cholis:2014lta,Yuan:2014rca,Petrovic:2014xra,Brandt:2015ula,Malyshev:2024obk,Kuvatova:2024bdn}. This hypothesis is motivated by the fact that pulsars produce gamma-ray emission with a spectral shape that is similar to that of the GCE~\cite{Baltz:2006sv,Hooper:2010mq,Hooper:2011ti,Abazajian:2012pn,Gordon:2013vta}. 

Pulsars are rapidly spinning neutron stars which steadily transfer their rotational kinetic energy into emission at radio, x-ray, and gamma-ray wavelengths. Most pulsars are the product of supernovae, which leave behind a neutron star with a large surface magnetic field ($B\sim 10^{11}-10^{13} \, \textrm{G}$) and an initial rotational period on the order of $P \sim 0.1 \, \textrm{s}$. Magnetic dipole breaking causes a neutron star to lose its rotational kinetic energy on the following timescale:
%
\begin{align}
\tau \equiv \frac{E}{\dot{E}} = \frac{P}{2\dot{P}} \approx 1.6 \times 10^5 \, \textrm{yr} \times \bigg(\frac{10^{12} \, \textrm{G}}{B}\bigg)^2 \bigg(\frac{P}{0.1 \, \textrm{s}}\bigg)^2.
\end{align}
%
Due to their large magnetic fields, young pulsars lose rotational kinetic energy and become faint relatively quickly. In contrast, some neutron stars are ``spun-up'' to high rates of rotation through dynamical interactions with a stellar companion. These ``recycled'' pulsars typically feature periods on the order of $P\sim 2-10$ ms and magnetic fields of $B\sim 10^8-10^9 \, \textrm{G}$, allowing them to retain their rotational kinetic energy over much longer timescales, \mbox{$\tau \sim 0.1-100 \, \textrm{Gyr}$~\cite{1982Natur.300..728A,Lorimer:2001vd,Phinney:1994gf,Lorimer:2008se,Kiziltan:2009rx}.} The long lifetimes of MSPs make it plausible that large numbers of these objects could be present in the Inner Galaxy, having been produced in situ, or through the dynamical friction and tidal disruption of globular clusters (which are known to contain large populations of MSPs)~\cite{Brandt:2015ula,Fragione:2017rsp,Gnedin:2013cda}. 


Pulsar interpretations of the GCE were elevated substantially in 2015, when two independent groups claimed to have identified evidence that the gamma rays associated with the GCE are spatially clustered, as might be expected if this signal is produced by a population of near-threshold point sources~\cite{Lee:2015fea,Bartels:2015aea}. These early analyses, however, have not since held up to scrutiny~\cite{Leane:2019xiy,Leane:2020nmi,Leane:2020pfc,Zhong:2019ycb}. In particular, it has been shown that the non-Poissonian template fit employed in Ref.~\cite{Lee:2015fea} tends to misattribute smooth gamma-ray signals to point source populations, resulting from imperfections in the modeling of the diffuse backgrounds~\cite{Leane:2019xiy,Leane:2020nmi,Leane:2020pfc}, though the effect of these systematic errors on the preference of such models for pulsar interpretations is disputed by Refs.~\cite{Chang:2019ars, Buschmann:2020adf}.  Furthermore, it was demonstrated in Ref.~\cite{Zhong:2019ycb} that when updated point source catalogs are utilized, wavelet-based analysis (such as those employed in Ref.~\cite{Bartels:2015aea}) do
not favor point source interpretations of the GCE. In light of these results, there is now general agreement that the GCE could arise from either a smooth distribution of sources, such as from annihilating dark matter, or from a very large number of faint point sources (for more recent attempts to address this question using machine learning techniques, see Refs.~\cite{List:2021aer,Mishra-Sharma:2021oxe,List:2020mzd,Caron:2017udl, Caron:2022akb, Ramirez:2024oiw}). 

Pulsar interpretations of the GCE were also bolstered by claims that the spatial morphology of the GCE traces known stellar populations, such as the Galactic Bulge and Bar~\cite{Macias:2019omb,Macias:2016nev,Bartels:2017vsx,Ramirez:2024oiw,Song:2024iup}. If confirmed, this would favor an astrophysical origin for this signal. Once again, significant doubt has since been cast on these claims. In particular, several studies have shown that either dark matter-like or bulge-like morphologies can provide a better fit to the gamma-ray data, depending on the background model and masking procedure that is adopted~\cite{Zhong:2024vyi,McDermott:2022zmq,Cholis:2021rpp,DiMauro:2021raz, Song:2024iup}. 


Several arguments have been leveled against 
MSP interpretations of the GCE:
\begin{enumerate}
%
\item{Low-mass X-ray binaries (LMXBs) are the primary progenitors of MSPs~\cite{1982Natur.300..728A,2009Sci...324.1411A}, and yet relatively few of these objects have been observed in the direction of the Inner Galaxy~\cite{Cholis:2014lta}. By comparing the number of LMXBs in the Inner Galaxy to those found in globular clusters, and scaling that ratio to the gamma-ray emission observed from globular clusters, it can estimated that only 4-11\% of the gamma-ray emission associated with the GCE can be attributed to MSPs~\cite{Haggard:2017lyq}.}
%
\item{Pulsars typically receive natal velocity kicks at the time of their formation, typically on the order $\sim 10^2~\textrm{km/s}$. These kicks should cause the spatial distribution of the resulting MSPs to be broader than that of the overall stellar population, and inconsistent with the observed morphology of the GCE, especially regarding the intensity of the excess observed within the innermost $\sim 1^{\circ}$ around the Galactic Center~\cite{Boodram:2022lpn}.}
%
\item{Very few pulsars have been detected in the vicinity of the Inner Galaxy. This allows us to place strong constraints on the gamma-ray luminosity function of any MSP population that could be responsible for the GCE. These constraints indicate that if MSPs generate the gamma-ray excess, those pulsars must be systematically less bright than those found in other environments. This, in turn, requires a larger number of MSPs to generate the observed intensity of the GCE, on the order of $N_\textrm{GCE, tot} \gsim 10^5$~\cite{Holst:2024fvb} (see also, Refs.~\cite{Cholis:2014noa,Hooper:2016rap,Hooper:2015jlu,Bartels:2018xom,Ploeg:2017vai,Ploeg:2020jeh,Dinsmore:2021nip}).}
%
\end{enumerate}

It has been suggested that these challenges could be mitigated if a large fraction of the Inner Galaxy's MSP population is not formed through the transfer of angular momentum to neutron stars, but instead through the accretion-induced collapse of O-Ne white dwarfs~\cite{Gautam:2021wqn}. Unlike ordinary MSPs, this hypothetical pulsar population would not be evolutionary connected to LMXBs, and their members would not be expected to receive large natal kicks, allowing them to evade the first two of the arguments listed above. It is not at all clear, however, why a large fraction of MSPs in the Inner Galaxy would be produced through accretion induced collapse when this is not the case among pulsar populations observed in other environments. 
% 
In fact, it has generally
been found that accretion-induced collapse only rarely produces binary configurations that are amenable to further mass-transfer-based spin-up of the compact object \cite{Kalogera_1996, Willems_2004} and moreover produces neutron stars with magnetic fields too
small to produce MSP activity \cite{Wijers_1997, sutantyo2000formationbinarymillisecondpulsars}.

%
Alternatively, if the magnetic fields of MSPs formed through accretion-induced collapse are of the same magnitude as MSPs formed through standard means, there would be no reason to expect either of these two populations to be systematically fainter than the other.



The population of MSPs observed in the Galactic Plane contains many pulsars that would have been detected by Fermi if they had been present in the Inner Galaxy. In particular, if the GCE is produced by MSPs with the same luminosity function and other characteristics as those observed within the Galactic Plane, $N_\textrm{MSP} \sim 10-35$ of the pulsars associated with the Inner Galaxy population should have been been detected by Fermi and included within their 4FGL and 3PC source catalogs~\cite{Holst:2024fvb,Fermi-LAT:2023zzt}. As only three such candidate pulsars exist~\cite{Holst:2024fvb}, such a scenario is in significant tension with this possibility.

It remains possible, however, that the Inner Galaxy could contain a large number of MSPs that are systematically fainter or otherwise more difficult to identify than those observed within the Galactic Plane. Motivating this possibility is the fact that the Galactic Bulge experienced a high rate of star formation when it was young, allowing for the formation of a large number of MSPs which are now very old, potentially causing them be systematically lower in luminosity than those found in the Galactic Plane. The stellar populations present within globular clusters are also very old, however, comparable in age to that of the Galactic Bulge. It is important, therefore, to measure the gamma-ray luminosity function of MSPs in globular clusters, as this could have great bearing on the viability of pulsar interpretations of the GCE. 

In this study, we revisit our previous work~\cite{Hooper:2016rap} attempting to constrain the gamma-ray luminosity function of MSPs in globular clusters. To this end, we use the publicly available Fermi data to measure or constrain the gamma-ray emission from each of the Milky Way's 157 globular clusters (for related work, see Refs.~\cite{Hou:2024wpr, ballet2024fermilargeareatelescope,Zhang:2023gwr,Evans:2022zno,Zhang:2022jlx,Yuan:2022egs,deMenezes:2018ilq,Cholis:2014noa,Fermi-LAT:2010gpc}). We detect statistically significant gamma-ray emission from 56 of these globular clusters, 8 of which are not contained in previous
gamma-ray source catalogs. We then use this data, in combination with the visible luminosities and previously calculated stellar encounter rates of the globular clusters, to constrain the gamma-ray luminosity function of the MSPs in these systems. We find that this MSP population has a luminosity function that is consistent with that previously reported for MSPs in the Galactic Plane, featuring a mean luminosity of $\langle L_{\gamma}\rangle \sim (1-8) \times 10^{33} \, \textrm{erg/s}$. If the GCE is generated by pulsars with the same gamma-ray luminosity function and other characteristics as those found in globular clusters, Fermi's source catalogs would have contained $N_\textrm{MSP} \sim 17-37$ pulsars associated with this population. From this, we conclude that the GCE could be generated by pulsars only if they are significantly lower in luminosity or otherwise more difficult to identify than those in globular clusters (or in the Galactic Plane). 

